\documentclass[11pt]{article}
\usepackage[margin=1in]{geometry}
\usepackage{listings}
\usepackage{xcolor}
\usepackage{hyperref}
\usepackage{parskip}

\lstdefinestyle{python}{
    language=Python,
    basicstyle=\ttfamily\small,
    keywordstyle=\color{blue},
    commentstyle=\color{gray},
    stringstyle=\color{orange},
    showstringspaces=false,
    frame=single,
    breaklines=true
}

\title{Python Engineering Notes}
\author{}
\date{}

\begin{document}
\maketitle
\tableofcontents
\newpage

% -------------------------------------------------------
\section{Object-Oriented Programming}

Python is an object-oriented language in which a \textit{class} serves as a blueprint
for objects that bundle data (attributes) and behaviour (methods) together.
A class is defined with the \texttt{class} keyword, and its constructor is the
special method \texttt{\_\_init\_\_}, which is called automatically whenever a new
instance is created.

\begin{lstlisting}[style=python]
class RiskMeasures:
    def __init__(self, samples, alpha):
        self.samples = samples
        self.alpha   = alpha
        self.evar    = 0.0
\end{lstlisting}

Here \texttt{self} refers to the instance being constructed. Every method receives
\texttt{self} as its first argument so that it can read and modify the instance's
own attributes.

\subsection{The \texttt{pass} statement}

\texttt{pass} is a syntactic no-op: it tells Python ``there is nothing to do here.''
Python requires every indented block to contain at least one statement, so
\texttt{pass} exists purely to satisfy that requirement when a block would otherwise
be empty.

It is \textbf{necessary} in two common situations. First, when writing a class or
function whose body has not been implemented yet:

\begin{lstlisting}[style=python]
class NotYetImplemented:
    pass          # required: class body cannot be empty

def placeholder():
    pass          # required: function body cannot be empty
\end{lstlisting}

Second, when you intentionally want to suppress an exception without taking any
action:

\begin{lstlisting}[style=python]
try:
    result = 1 / x
except ZeroDivisionError:
    pass          # required: except block cannot be empty
\end{lstlisting}

It is \textbf{useless} (and misleading) as soon as any real code exists in the
block. A \texttt{pass} that follows a \texttt{return} statement is unreachable,
because the function has already exited:

\begin{lstlisting}[style=python]
def CVaR(self):
    self.var  = np.percentile(self.samples, 100 * self.alpha)
    C         = self.samples[self.samples >= self.var]
    self.cvar = np.mean(C)
    return self.cvar
    pass          # never reached; remove this
\end{lstlisting}

Similarly, a \texttt{pass} inside \texttt{\_\_init\_\_} after real assignments does
nothing and should be removed to keep the code clean.

The rule of thumb is simple: use \texttt{pass} only when the block has
\textit{nothing else} in it. The moment you write any real statement, \texttt{pass}
becomes noise.

\subsection{Consistent state management in methods}

When a class stores computed values as instance attributes, every method that
produces such a value should update the corresponding attribute before returning it.
This keeps the object's state consistent and allows other methods to rely on it
without recomputing. In the \texttt{RiskMeasures} class, for example, \texttt{EVaR}
both sets \texttt{self.evar} and returns the value, so that either usage pattern
works correctly:

\begin{lstlisting}[style=python]
def EVaR(self):
    def f(z, alp, s):
        c       = z * np.max(s)
        log_mgf = c + np.log(np.mean(np.exp(z * s - c)))
        return (1 / z) * (log_mgf - np.log(1 - alp))

    res        = minimize_scalar(f, bounds=(1e-6, 100),
                                 method='bounded',
                                 args=(self.alpha, self.samples))
    self.evar  = res.fun
    return res.fun
\end{lstlisting}

% -------------------------------------------------------
\section{Dictionaries}

A dictionary maps keys to values. Keys are typically strings or integers; values
can be anything, including NumPy arrays.

\begin{lstlisting}[style=python]
distributions = {
    'Gaussian':   np.random.normal(0, 1, 1000),
    'Log-normal': np.random.lognormal(0, 1, 1000),
}
\end{lstlisting}

\subsection{Iterating over a dictionary}

Looping over a dictionary directly gives you only the \textbf{keys}:

\begin{lstlisting}[style=python]
for i in distributions:
    print(i)        # prints 'Gaussian', 'Log-normal', ...
    print(i.value)  # AttributeError: i is a string, not an object with .value
\end{lstlisting}

To get both the key and the value simultaneously, use \texttt{.items()}, which
returns each entry as a \texttt{(key, value)} pair:

\begin{lstlisting}[style=python]
for name, samples in distributions.items():
    print(name)     # 'Gaussian', 'Log-normal', ...
    print(samples)  # the numpy array
\end{lstlisting}

You can also iterate over only the values with \texttt{.values()}, or only the
keys with \texttt{.keys()}:

\begin{lstlisting}[style=python]
for samples in distributions.values():  # values only
    ...

for name in distributions.keys():       # keys only (same as plain loop)
    ...
\end{lstlisting}

In practice, \texttt{.items()} is the most common pattern when both the key and
the value are needed --- as when looping over distributions and using the name as
a plot title.

\subsection{\texttt{enumerate} --- loop with an index}

When you need both a counter and the values from an iterable, use
\texttt{enumerate} instead of maintaining a manual counter variable:

\begin{lstlisting}[style=python]
# fragile: i resets or increments in the wrong place
i = 0
for name, samples in distributions.items():
    ...
    i += 1

# clean: enumerate gives (index, value) pairs automatically
for i, (name, samples) in enumerate(distributions.items()):
    ...
\end{lstlisting}

\texttt{enumerate} starts counting from 0 by default. The start value can be
changed: \texttt{enumerate(items, start=1)}.

% -------------------------------------------------------
\section{Matplotlib subplots}

\texttt{plt.subplots(2, 2)} creates a $2\times2$ grid of axes and returns them
as a 2D NumPy array of shape \texttt{(2, 2)}. Indexing it with a single integer
(\texttt{ax[0]}) returns a whole row, not a single subplot. To index subplots
with a flat counter, call \texttt{.flatten()} first:

\begin{lstlisting}[style=python]
fig, ax = plt.subplots(2, 2, figsize=(10, 8))
axes = ax.flatten()   # shape (4,): axes[0] ... axes[3]

for i, (name, samples) in enumerate(distributions.items()):
    axes[i].plot(alpha_range, hist_c, label='CVaR')
    axes[i].plot(alpha_range, hist_e, label='EVaR')
    axes[i].set_title(f'{name}')
    axes[i].set_xlabel('Alpha')
    axes[i].set_ylabel('Risk Measure')
    axes[i].legend()
    axes[i].grid(True)

plt.tight_layout()
plt.show()
\end{lstlisting}

With subplots, labels and titles are set on each axis object
(\texttt{axes[i].set\_xlabel(...)}) rather than at the figure level
(\texttt{plt.xlabel(...)}). \texttt{plt.tight\_layout()} prevents titles and
labels from overlapping across subplots.

\end{document}
